% Part:sets-functions-relations
% Chapter: size-of-sets
% Section: zig-zag

\documentclass[../../../include/open-logic-section]{subfiles}

\begin{document}

\olfileid[ar]{sfr}{siz}{zigzag}

\olsection{طريقة كانتور المتعرجة}

\begin{explain}
تقدمت تعدادات يسيرة، فلننتقل إلى ما هو أدق قليلًا: مجموعة أزواج
الطبيعيّات\oliflabeldef{sfr:set:pai:sec}{، التي سبق تعريفها في
\olref[set][pai]{sec} بهذه الصورة:}{ التي تعريفها:}
\[
\Nat \times \Nat = \Setabs{\tuple{n,m}}{n,m \in \Nat}
\]
ولنجعل هذه الأزواج في \emph{مصفوفة}:
\[
\begin{array}{ c | c | c | c | c | c}
& \mathbf 0 & \mathbf 1 & \mathbf 2 & \mathbf 3 & \dots \\
\hline
\mathbf 0 & \tuple{0,0} & \tuple{0,1} & \tuple{0,2} & \tuple{0,3} & \dots \\
\hline
\mathbf 1 & \tuple{1,0} & \tuple{1,1} & \tuple{1,2} & \tuple{1,3} & \dots \\
\hline
\mathbf 2 & \tuple{2,0} & \tuple{2,1} & \tuple{2,2} & \tuple{2,3} & \dots \\
\hline
\mathbf 3 & \tuple{3,0} & \tuple{3,1} & \tuple{3,2} & \tuple{3,3} & \dots \\
\hline
\vdots & \vdots & \vdots & \vdots & \vdots & \ddots\\
\end{array}
\]
فكل زوج في $\Nat \times \Nat$ له في المصفوفة موضع واحد لا غير؛
فالزوج $\tuple{n,m}$ في الصف $n$ والعمود $m$. فكيف نحول صفوفها
وأعمدتها إلى قائمة ذات بعد واحد؟ لهذا وجوه كثيرة، ومنها نمط المصفوفة الآتية:
\[
\begin{array}{ c | c | c | c | c | c | c}
& \mathbf 0 & \mathbf 1 & \mathbf 2 & \mathbf 3 & \mathbf 4 &\dots \\
\hline
\mathbf 0 & 0  & 1& 3 & 6& 10 &\ldots \\
\hline
\mathbf 1 &2 & 4& 7 & 11 & \dots &\ldots \\
\hline
\mathbf 2 & 5 & 8 & 12 & \ldots & \dots&\ldots \\
\hline
\mathbf 3 & 9 & 13 & \ldots & \ldots & \dots & \ldots \\
\hline
\mathbf 4 & 14 & \ldots & \ldots & \ldots & \dots & \ldots \\
\hline
\vdots & \vdots & \vdots & \vdots & \vdots&\ldots & \ddots\\
\end{array}
\]\noindent
واسم هذا النمط \emph{طريقة كانتور المتعرجة}، وبه يكون تعداد
$\Nat \times \Nat$:
\[
\tuple{0,0}, \tuple{0,1}, \tuple{1,0}, \tuple{0,2}, \tuple{1,1},
\tuple{2,0}, \tuple{0,3}, \tuple{1,2}, \tuple{2,1}, \tuple{3,0}, \dots
\]
فقد ثبت الآتي:
\end{explain}

\begin{prop}\ollabel{natsquaredenumerable}
المجموعة $\Nat \times \Nat$ هي !!{enumerable}.
\end{prop}

\begin{proof}
عرّف $f \colon \Nat \to \Nat\times\Nat$ بجعل صورة $k \in \Nat$
الزوج $\tuple{n,m} \in \Nat \times \Nat$ الذي يقع فيه $k$
في الصف $n$ والعمود $m$ من مصفوفة كانتور المتعرجة.
\end{proof}

\begin{explain}
ولهذه الطريقة تعميم يسير. فمنه تعداد الثلاثيات المرتبة من
الطبيعيّات، أي:
\[
\Nat \times \Nat \times \Nat = \Setabs{\tuple{n,m,k}}{n,m,k \in \Nat}
\]
نأخذ $\Nat \times \Nat \times \Nat$ ضربًا ديكارتيًا
لـ$\Nat \times \Nat$ في $\Nat$، أي
\[
\Nat^3 = (\Nat \times \Nat) \times \Nat =
\Setabs{\tuple{\tuple{n,m},k}}{n, m, k
  \in \Nat }
\]
فنعدّد $\Nat^3$ بمصفوفة أحد محوريها تعداد $\Nat$،
والآخر تعداد $\Nat^2$:
\[
\begin{array}{ c | c | c | c | c | c}
& \mathbf 0 & \mathbf 1 & \mathbf 2 & \mathbf 3 & \dots \\
\hline
\mathbf{\tuple{0,0}} & \tuple{0,0,0} & \tuple{0,0,1} & \tuple{0,0,2} & \tuple{0,0,3} & \dots \\
\hline
\mathbf{\tuple{0,1}} & \tuple{0,1,0} & \tuple{0,1,1} & \tuple{0,1,2} & \tuple{0,1,3} & \dots \\
\hline
\mathbf{\tuple{1,0}} & \tuple{1,0,0} & \tuple{1,0,1} & \tuple{1,0,2} & \tuple{1,0,3} & \dots \\
\hline
\mathbf{\tuple{0,2}} & \tuple{0,2,0} & \tuple{0,2,1} & \tuple{0,2,2} & \tuple{0,2,3} & \dots\\
\hline
\vdots & \vdots & \vdots & \vdots & \vdots & \ddots \\
\end{array}
\]
ثم نسلك نظير طريق كانتور المتعرج فنحصل على تعداد~$\Nat^3$.
وبتكرار الصنعة نحصل على تعداد $\Nat^n$ لأي طبيعي $n$، فيثبت:
\end{explain}

\begin{prop}
المجموعة $\Nat^n$ هي !!{enumerable} لكل $n \in \Nat$.
\end{prop}

\begin{prob}\label{sfr:siz:zigzag:prob:posint-n}
برهن على أن $(\PosInt)^n$ هي !!{enumerable} لكل $n \in \Nat$.
\end{prob}

\begin{prob}\label{sfr:siz:zigzag:prob:posint-star}
  برهن على أن $(\PosInt)^*$ هي !!{enumerable}، ولك استعمال \cref{sfr:siz:zigzag:prob:posint-n}.
\end{prob}

\end{document}
